\begin{table}[!htbp]
\centering\small
\caption{H2: korelacja Spearmana między średnią wartością miary a~wynikiem zespołu po 15 wariantach, liczona w~każdym zbiorze osobno (bagging, klasyfikacja). Podano medianę po zbiorach, liczbę zbiorów z~dodatnią korelacją i~wartość $p_H$ testu Wilcoxona z~korektą Holma (7 miar). Korelacja cząstkowa uwzględnia średnią jakość drzew.}
\label{tab:h2bag}
\begin{tabular}{lrrrrrr}
\toprule
& \multicolumn{3}{c}{\textbf{korelacja}} & \multicolumn{3}{c}{\textbf{cząstkowa}} \\
\cmidrule(lr){2-4}\cmidrule(lr){5-7}
\textbf{Miara} & $\tilde\rho$ & $\rho>0$ & $p_H$ & $\tilde\rho$ & $\rho>0$ & $p_H$ \\
\midrule
\multicolumn{7}{l}{\emph{Klasyfikacja}}\\
$D_J$ & -0.16 & 5/12 & 1.000 & +0.54 & 12/12 & \textbf{0.003} \\
$D_C$ & -0.09 & 5/12 & 1.000 & +0.60 & 12/12 & \textbf{0.003} \\
$D_Y$ & -0.24 & 5/12 & 1.000 & +0.54 & 12/12 & \textbf{0.003} \\
$U$ & +0.50 & 8/12 & 1.000 & +0.12 & 8/12 & 0.339 \\
$R_Y$ & -0.36 & 4/12 & 1.000 & +0.26 & 9/12 & 0.068 \\
$R_X$ & -0.37 & 5/12 & 1.000 & +0.41 & 10/12 & \textbf{0.028} \\
różnorodność predykcji & -0.16 & 5/12 & 1.000 & +0.62 & 12/12 & \textbf{0.003} \\
\bottomrule
\end{tabular}
\end{table}
